% Source printed pp. 144-148, physical PDF pp. 152-156.
\section{光滑性与无穷小提升判据}\label{src:4.7-section}
\begin{epigraph}
真理总能在简洁之中找到, 而不在事物的繁多与混乱之中.
\epigraphauthor{Isaac Newton}
\end{epigraph}
概形的\newterm{无穷小提升判据}{infinitesimal lifting criterion} (\bookref{A.3.1}(3), \bookref{A.3.2}(5) 和 \bookref{A.3.3}(4)) 从函子的角度刻画光滑, 平展或非分歧态射. 下面将其推广到代数叠 (定理~\bookref{4.7.1}). 作为应用, 我们证明光滑曲线的模叠 $\mathcal M_g$ 及固定光滑曲线上凝聚层的模叠 $\Coh_{r,d}(C)$ 都光滑 (命题~\bookref{4.7.7} 和~\bookref{4.7.9}).
\subsection{无穷小提升判据}\label{src:4.7.1-subsection}
回忆, 光滑性是源和靶上的光滑局部性质, 因而对代数叠的态射而言, 光滑性是良定义的. 另一方面, 平展性和非分歧性在靶上是光滑局部的, 但在源上只是平展局部的, 因而对相对 Deligne--Mumford 的态射 $\mathcal X\to\mathcal Y$ 而言是良定义的. 这里相对 Deligne--Mumford 是指, 对每个从概形出发的映射 $T\to\mathcal Y$, 基变换 $\mathcal X\times_{\mathcal Y}T$ 都是 Deligne--Mumford 的; 见定义~\bookref{4.3.36}. 代数叠的态射 $\mathcal X\to\mathcal Y$ 是平展的 (相应地, 非分歧的), 若它是相对 Deligne--Mumford 的, 且每个基变换 $\mathcal X\times_{\mathcal Y}T\to T$ 都平展 (相应地, 非分歧); 见定义~\bookref{4.3.37}. 还需回忆, $\mathcal X\to\mathcal Y$ 相对 Deligne--Mumford 当且仅当对角线 $\mathcal X\to\mathcal X\times_{\mathcal Y}\mathcal X$ 非分歧 (推论~\bookref{4.6.5}).
\begin{theorem}[光滑性/平展性/非分歧性的无穷小提升判据]\label{src:4.7.1}
设 $f:\mathcal X\to\mathcal Y$ 是局部 Noether 代数叠之间的局部有限型态射, 两个叠的对角线都拟紧且分离. 考虑由实线箭头构成的 $2$-交换图
\begin{equation}\label{src:4.7.2}
\begin{tikzcd}
\Spec A_0\arrow[r]\arrow[d,hook]&\mathcal X\arrow[d,"f"]\\
\Spec A\arrow[r]\arrow[ur,dashed]&\mathcal Y,
\end{tikzcd}
\end{equation}
其中 $A\to A_0$ 是剩余域为 $k$ 的 Artin 局部环之间的满射, 且 $\ker(A\to A_0)\cong k$. 则
\begin{enumerate}[(1)]
\item $f$ 光滑当且仅当每个图~\bookeqref{4.7.2} 都存在提升;
\item $f$ 平展当且仅当每个图~\bookeqref{4.7.2} 都存在提升, 且提升在唯一同构的意义下唯一;
\item $f$ 非分歧当且仅当图~\bookeqref{4.7.2} 的任意两个提升都唯一同构;
\item $f$ 的对角线非分歧当且仅当图~\bookeqref{4.7.2} 的提升的每个自同构都平凡.
\end{enumerate}
\end{theorem}
\begin{remark}\label{src:4.7.3}
具体来说, $2$-交换图
\begin{equation}\label{src:4.7.4}
\begin{tikzcd}[column sep=3em,row sep=3em]
S\arrow[r,"x"]\arrow[d,"g"']&\mathcal X\arrow[d,"f"]\\
T\arrow[r,"y"']&\mathcal Y,
\arrow[Rightarrow,from=1-2,to=2-1,shorten <=1.7em,shorten >=1.7em,"\alpha"']
\end{tikzcd}
\end{equation}
的一个\newterm{提升}{lifting} 是三元组 $(\widetilde x,\beta,\gamma)$, 其中 $\widetilde x:T\to\mathcal X$ 是映射, $\beta:x\xrightarrow{\sim}\widetilde x\circ g$ 和 $\gamma:f\circ\widetilde x\xrightarrow{\sim}y$ 是 $2$-同构, 使下图交换:
\[
\begin{tikzcd}[column sep=small]
{}&f\circ x\arrow[dl,"f(\beta)"']\arrow[dr,"\alpha"]&{}\\
f\circ\widetilde x\circ g\arrow[rr,"g^*\gamma"']&{}&y\circ g
\end{tikzcd}
\]
可以将一个提升看作 $2$-交换图
\[
\begin{tikzcd}[column sep=4em,row sep=4em]
S\arrow[r,"x",""{name=top}]\arrow[d,"g"']&\mathcal X\arrow[d,"f",""{name=right,pos=.55}]\\
T\arrow[r,"y"',""{name=bottom}]\arrow[ur,"\widetilde x" description,""{name=upper,pos=.72},""{name=lower,pos=.28}]&\mathcal X.
\arrow[from=top,to=upper,Rightarrow,"\beta"']
\arrow[from=right,to=bottom,Rightarrow,shorten <=.4em,shorten >=.4em,"\gamma"']
\end{tikzcd}
\]
提升之间的态射 $(\widetilde x,\beta,\gamma)\to(\widetilde x',\beta',\gamma')$ 是 $2$-同构 $\Theta:\widetilde x\xrightarrow{\sim}\widetilde x'$, 满足 $\beta=g^*\Theta\circ\beta'$ 和 $\gamma=\gamma'\circ f(\Theta)$.

也可以用群胚之间的映射
\[
\Psi:\mathcal X(T)\longrightarrow\mathcal X(S)\times_{\mathcal Y(S)}\mathcal Y(T)
\]
解释提升. 图~\bookeqref{4.7.4} 定义对象 $(x,y,\alpha)\in\mathcal X(S)\times_{\mathcal Y(S)}\mathcal Y(T)$, 而提升范畴就是这个对象上的纤维范畴. 上述判据断言, $f:\mathcal X\to\mathcal Y$ 光滑 (相应地, 平展, 非分歧, 对角线非分歧), 当且仅当 $\mathcal X(A)\to\mathcal X(A_0)\times_{\mathcal Y(A_0)}\mathcal Y(A)$ 本质满 (相应地, 是等价, 完全忠实, 忠实).
\end{remark}
\begin{proof}
我们处理光滑性的判据, 其余情形留给读者. 假设 $f:\mathcal X\to\mathcal Y$ 光滑, 并给定图~\bookeqref{4.7.2}. 将 $\mathcal Y$ 替换为 $\Spec A$, 将 $\mathcal X$ 替换为 $\mathcal X\times_{\mathcal Y}\Spec A$ 后, 可假设 $\mathcal Y=\Spec A$ 仿射. 需要证明, $\Spec A_0$ 上的截面
\[
\begin{tikzcd}
{}&\mathcal X\arrow[d]&{}\\
\Spec A_0\arrow[r,hook]\arrow[ur]&\Spec A\arrow[u,dashed,bend right=40]&{}
\end{tikzcd}
\]
可延拓为 $\Spec A$ 上的截面. 若 $\mathcal X$ 是概形, 则概形的光滑性无穷小提升判据~\bookeqref{A.3.1} 给出提升的存在性. 若 $\mathcal X$ 是代数空间, 则可应用命题~\bookref{5.3.4} (将在后面独立证明), 构造从概形出发的光滑呈示 $U\to\mathcal X$, 以及 $\Spec k\hookrightarrow\Spec A_0\to\mathcal X$ 的提升 $\Spec k\to U$. 得到下图:
\[
\begin{tikzcd}[column sep=3em,row sep=2em]
{}&{}&U\arrow[d]\\
{}&{}&\mathcal X\arrow[d]\\
\Spec k\arrow[r,hook]\arrow[uurr]&\Spec A_0\arrow[r,hook]\arrow[ur]\arrow[uur,dashed]&\Spec A\arrow[uu,dashed,bend right=35].
\end{tikzcd}
\]
由于 $U\to\mathcal X$ 可由概形表出且光滑, 对平方零扩张 $A_0/\mathfrak m^{k+1}\to A_0/\mathfrak m^k$ 归纳地应用无穷小提升判据~\bookeqref{A.3.1}, 可得到 $\Spec A_0\to\mathcal X$ 的提升 $\Spec A_0\to U$. 由于 $U\to\Spec A$ 是概形的光滑态射, 再次应用无穷小提升判据, 可得到 $\Spec A_0\to U$ 的提升 $\Spec A\to U$, 而复合 $\Spec A\to U\to\mathcal X$ 就是所需的延拓. 最后, 若 $\mathcal X$ 是代数叠, 则可以重复上述论证, 应用命题~\bookref{5.3.4} 构造带有提升 $\Spec k\to U$ 的光滑呈示 $U\to\mathcal X$. 这里利用 $U\to\mathcal X$ 的可表性和代数空间的情形, 构造提升 $\Spec A_0\to U$.

反过来, 假设 $f:\mathcal X\to\mathcal Y$ 的每个图~\bookeqref{4.7.2} 都存在提升. 选取光滑呈示 $V\to\mathcal Y$ 和 $U\to\mathcal X_V$, 得到交换图
\[
\begin{tikzcd}
U\arrow[r]\arrow[dr]&\mathcal X_V\arrow[r]\arrow[d]&\mathcal X\arrow[d,"f"]\\
{}&V\arrow[r]\arrow[ur,phantom,"\square" description]&\mathcal Y.
\end{tikzcd}
\]
由于 $U\to\mathcal X_V$ 光滑, 由已经证明的方向, $U\to\mathcal X_V$ 的每个图~\bookeqref{4.7.2} 都存在提升. 由于 $\mathcal X_V\to V$ 的每个图~\bookeqref{4.7.2} 也存在提升 (因为它是 $\mathcal X\to\mathcal Y$ 的基变换), 复合 $U\to V$ 也具有这一性质. 由于 $U\to V$ 是概形的态射, 无穷小提升判据~\bookeqref{A.3.1} 蕴含它光滑. 因为光滑性是源和靶上的光滑局部性质, $\mathcal X\to\mathcal Y$ 光滑. 另见 \cite[命题 $4.15$]{LMB00} 和 \cite[Tag 0DP0]{SP}.
\end{proof}
\begin{exercise}\label{src:4.7.5}
证明无穷小提升判据中关于平展性, 非分歧性和对角线非分歧性的其余部分.

\textit{提示:} 对于对角线非分歧的情形, 将态射 $f:\mathcal X\to\mathcal Y$ 的提升的一个自同构, 与对角线 $\Delta_f:\mathcal X\to\mathcal X\times_{\mathcal Y}\mathcal X$ 的一个提升联系起来. 将平展性的情形归约到非分歧性. 对于非分歧性, 利用推论~\bookref{4.6.5}, 将问题归约为证明相对 Deligne--Mumford 态射 $f:\mathcal X\to\mathcal Y$ 非分歧当且仅当图~\bookeqref{4.7.2} 的任意两个提升都同构, 并用类似于光滑性判据证明的方法建立这一结论.
\end{exercise}
作为第一个应用, 我们看到极小呈示的存在性 (定理~\bookref{4.6.1}) 所构造的呈示是极小万有的, 即在所选原像处的切空间上诱导同构. 同时, 也可以用光滑代数叠的切空间与稳定子表示它的维数.
\begin{proposition}\label{src:4.7.6}
设 $\mathcal X$ 是 Noether 代数叠, $x\in|\mathcal X|$ 是稳定子光滑的有限型点. 设 $f:(U,u)\to(\mathcal X,x)$ 是从概形出发的光滑态射, 且 $\mathcal G_x\times_{\mathcal X}U\cong\Spec\kappa(u)$. 则 $U\to\mathcal X$ 在 $u$ 处极小万有, 即 $T_{U,u}\to T_{\mathcal X,f(u)}$ 是 $\kappa(u)$-向量空间同构. 特别地, 若 $\mathcal X$ 在域 $k$ 上光滑, 而 $x\in\mathcal X(L)$ 是稳定子光滑的点, 其中 $L/k$ 是有限扩张, 则
\[
\dim_x\mathcal X=\dim T_{\mathcal X,x}-\dim G_x.
\]
\end{proposition}
\begin{proof}
对光滑态射 $U\to\mathcal X$ 应用无穷小提升判据~\bookeqref{4.7.1}, 可得到 $T_{U,u}\to T_{\mathcal X,f(u)}$ 的满射性. 单射性由下图为 Cartesian 图这一事实得到:
\[
\begin{tikzcd}
\Spec\kappa(u)\arrow[r,hook]\arrow[d]&U\arrow[d]\\
\mathcal G_x\arrow[r,hook]\arrow[ur,phantom,"\square" description]&\mathcal X.
\end{tikzcd}
\]
事实上, 若 $\tau:\Spec\kappa(u)[\epsilon]\to U$ 是 $T_{U,u}$ 中映到 $0\in T_{\mathcal X,f(u)}$ 的元素, 则复合 $\Spec\kappa(u)[\epsilon]\to U\to\mathcal X$ 经过剩余 gerbe $\mathcal G_x$ 分解, 因而也经过 $\Spec\kappa(u)$ 分解. 所以 $\tau=0$. 对于最后的断言, 极小呈示的存在性~\bookeqref{4.6.1} 给出相对维数为 $\dim G_x$ 的光滑态射 $f:(U,u)\to(\mathcal X,x)$, 满足 $\mathcal G_x\times_{\mathcal X}U\cong\Spec\kappa(u)$. 因而 $\dim_x\mathcal X=\dim_u U-\dim G_x$. 又因为 $U$ 在 $u$ 处光滑, 有 $\dim_u U=\dim T_{U,u}=\dim T_{\mathcal X,x}$.
\end{proof}
\subsection{模叠的光滑性}\label{src:4.7.2-subsection}
将光滑性的无穷小提升判据~\bookeqref{4.7.1} 与形变理论结合, 既可以验证模叠的光滑性, 也可以计算其维数.
\begin{proposition}\label{src:4.7.7}
对于 $g\geq2$, Deligne--Mumford 叠 $\mathcal M_g$ 在 $\Spec\mathbb Z$ 上光滑, 相对维数为 $3g-3$.
\end{proposition}
\begin{proof}
设 $\Spec k\to\mathcal M_g$ 是从域 $k$ 出发的态射, 对应于 $k$ 上的光滑连通射影曲线 $C$. 考虑下图:
\begin{equation}\label{src:4.7.8}
\begin{tikzcd}[column sep=3em]
\Spec k\arrow[r]\arrow[rr,bend left=20,"{[C]}"]&\Spec A_0\arrow[r]\arrow[d,hook]&\mathcal M_g\arrow[d]\\
{}&\Spec A\arrow[r]\arrow[ur,dashed]&\Spec\mathbb Z,
\end{tikzcd}
\end{equation}
其中 $A\to A_0$ 是剩余域为 $k$ 的 Artin 局部环之间的满射, 且 $k=\ker(A\to A_0)$. 映射 $\Spec A_0\to\mathcal M_g$ 对应于曲线族 $C_0\to\Spec A_0$, 而图~\bookeqref{4.7.8} 中的提升对应于下列实线箭头图的延拓:
\[
\begin{tikzcd}
C\arrow[r,hook]\arrow[d]&C_0\arrow[r,hook,dashed]\arrow[d]&\mathcal C\arrow[d,dashed]\\
\Spec k\arrow[r,hook]&\Spec A_0\arrow[r,hook]&\Spec A.
\end{tikzcd}
\]
由形变理论~\bookeqref{C.2.4}, 存在上同调类 $\operatorname{ob}_C\in H^2(C,T_C)$, 使上图存在延拓当且仅当 $\operatorname{ob}_C=0$. 由于 $C$ 是曲线, $H^2(C,T_C)=0$, 因而 $\operatorname{ob}_C=0$, 存在提升. 最后, 形变理论~\bookeqref{C.2.4} 还给出等同 $T_{\mathcal M_g,[C]}=H^1(C,T_C)$, 而 Riemann--Roch 计算表明它的维数为 $3g-3$ (见例~\bookref{4.5.8}). 由于 $\dim\Aut(C)=0$, 由命题~\bookref{4.7.6} 得到 $\dim_{[C]}(\mathcal M_g\times_{\mathbb Z}k)=3g-3$.
\end{proof}
\begin{proposition}\label{src:4.7.9}
对于代数闭域 $k$ 上的光滑连通射影曲线 $C$, 代数叠 $\Coh_{r,d}(C)$ 在 $\Spec k$ 上光滑, 维数为 \mbox{$r^2(g-1)$}.
\end{proposition}
\begin{proof}
由习题~\bookref{4.7.10}, $\Coh_{r,d}(C)$ 非空. 设 $[F]\in\Coh_{r,d}(C)(L)$ 是 $C_L=C\times_k L$ 上秩为 $r$, 次数为 $d$ 的凝聚层. 设 $A\to A_0$ 是剩余域为 $L$ 的 Artin 局部环之间的满射, 且 $L=\ker(A\to A_0)$. 需要检验, $C_{A_0}$ 上每个限制为 $F$ 的凝聚层 $\mathcal F_0$, 都可以延拓为 $C_A$ 上的凝聚层 $\mathcal F$. 由形变理论~\bookeqref{C.2.11}, 存在元素 $\operatorname{ob}_F\in\operatorname{Ext}^2_{\mathcal O_C}(F,F)$, 使 $\operatorname{ob}_F=0$ 当且仅当存在延拓. 由于 $C$ 是光滑曲线, 有 $\operatorname{Ext}^2_{\mathcal O_C}(F,F)=H^2(C_L,F\otimes F^\vee)=0$. 形变理论还给出等同 $T_{\Coh_{r,d}(C),[F]}=\operatorname{Ext}^1_{\mathcal O_C}(F,F)$, 而 Riemann--Roch 计算给出 $\dim\operatorname{Ext}^1_{\mathcal O_C}(F,F)=\dim\Aut(F)+r^2(g-1)$ (见例~\bookref{4.5.9}). 因而
\[
\begin{aligned}
\dim_{[F]}\Coh_{r,d}(C)&=\dim\operatorname{Ext}^1_{\mathcal O_C}(F,F)-\dim\Aut(F)\\
&=r^2(g-1).
\end{aligned}
\]
\end{proof}
\begin{exercise}\label{src:4.7.10}
对于每个线丛 $L$ 和秩 $r>0$, 证明存在秩为 $r$ 的向量丛 $E$, 使 $\det E\cong L$.
\end{exercise}
